\documentclass{article}
\usepackage[utf8]{inputenc}
\usepackage[spanish]{babel}
\usepackage{enumitem}
\usepackage{framed}
\usepackage{hyperref}

\title{Desarrollo del Test}
\author{}
\date{}

\begin{document}

\maketitle

\section{Desarrollo del test}

\subsection{¿Por qué necesitamos una metodología?}

Lo necesitamos para una innovación, para tener ideas claras o más seguras de algún proyecto, o idea ya sea de negocios o nuestra vida cotidiana.

\begin{framed}
\noindent\textbf{Modificaciones sugeridas del documento:}
\begin{enumerate}
    \item Contexto histórico: Obsolescencia del code \& fix en los años 70.
    \item Problemas sin metodología: Requerimientos ambiguos y sin normas.
    \item Ventajas: División en etapas y definición de entradas/salidas.
    \item Objetivo: Prolijidad, corrección y control en cada etapa.
    \item Resultado: Aumenta las posibilidades de éxito del proyecto.
\end{enumerate}
\end{framed}

\subsection{¿Sirve el modelo de ciclo de vida Code \& Fix?}

Code \& fix, cómo tal lo dice el texto es codificar y corregir, podemos codificar nueva información para nuestra vida cotidiana, alimentarnos de mejores hábitos, tener una cálida de vids mejor, y ello nos sirve para corregir nuestros propios errores, claro está que no se podría llamar de tal forma si no también enseñanza. Estás nos quedas de enseñanza para no volver a cometer.

\begin{framed}
\noindent\textbf{Modificaciones sugeridas del documento:}
\begin{enumerate}
    \item Definición: Codificar y corregir sobre la marcha.
    \item Uso adecuado: Proyectos pequeños de 1 o 2 programadores.
    \item Uso inadecuado: Sistemas complejos o grandes.
    \item Desventajas: Costos mayores, tiempo excesivo y calidad dudosa.
    \item Riesgo: Pérdida de control sobre la gestión del proyecto.
\end{enumerate}
\end{framed}

\subsection{¿Existe algún modelo de ciclo de vida que predomine?}

Yo predominio una forma extrovertida, amable, amigable, sociales. Etc. Ya es que lo que más sobresale de mi persona.

\begin{framed}
\noindent\textbf{Modificaciones sugeridas del documento:}
\begin{enumerate}
    \item Predominio: Ningún modelo predomina sobre los otros.
    \item Selección: Elegir el que mejor se adapte al proyecto.
    \item Factores: Complejidad, tiempo y entregas parciales.
    \item Comunicación: Entre equipo de desarrollo y usuario.
    \item Consejo: La mejor metodología es la que mejor conocemos.
\end{enumerate}
\end{framed}

\subsection{¿Seguir un modelo de ciclo de vida, nos garantiza el éxito del desarrollo?}

si, si tenemos buenos hábitos, si tenemos buenos proyectos, varias indicaciones podremos llegar a tener una vida llena de éxitos, y una gran cálida de vida.

\begin{framed}
\noindent\textbf{Modificaciones sugeridas del documento:}
\begin{enumerate}
    \item Garantía: Ningún modelo evita todos los riesgos.
    \item Incertidumbre: No eliminan la incertidumbre del cambio.
    \item Éxito: No garantizan el éxito del desarrollo.
    \item Función: Ayudan a prepararse para cambios y problemas.
    \item Definición: El éxito se define en las etapas iniciales.
\end{enumerate}
\end{framed}

\subsection{¿Se puede medir la incertidumbre que tenemos sobre los requerimientos iniciales?}

No porque es algo mental y no podemos ver la gravedad del problema.

\begin{framed}
\noindent\textbf{Modificaciones sugeridas del documento:}
\begin{enumerate}
    \item Medición: Se puede medir como cantidad de información necesaria.
    \item Baja incertidumbre: Cascada puro, Lineal, En V.
    \item Alta incertidumbre: Espiral, Evolutivo, Prototipos.
    \item Fuente: A menudo es el propio cliente o usuario.
    \item Conocimiento: El usuario no sabe perfectamente las funcionalidades.
\end{enumerate}
\end{framed}

\subsection{¿La generación de programas prototipo, es exclusiva de un solo modelo ciclo de vida?}

\begin{framed}
\noindent\textbf{Modificaciones sugeridas del documento:}
\begin{enumerate}
    \item El uso de programas prototipo no es exclusivo del ciclo de vida iterativo.
    \item En la práctica los prototipos se utilizan para validar los requerimientos de los usuarios en cualquier ciclo de vida.
    \item Modelos que utilizan prototipos formalmente: Ciclo de Vida por Prototipos, En Espiral e Iterativo.
    \item Se recomienda cuando no se conocen especificaciones precisas o hay innovaciones importantes.
    \item Según el documento (página 28): "Si no se conoce exactamente cómo desarrollar un determinado producto o cuáles son las especificaciones de forma precisa, suele recurrirse a definir especificaciones iniciales para hacer un prototipo".
\end{enumerate}
\end{framed}

\subsection{¿Podemos utilizar un lenguaje imperativo para el modelo de ciclo de vida orientado a objetos?}

\begin{framed}
\noindent\textbf{Modificaciones sugeridas del documento:}
\begin{enumerate}
    \item Sí se puede usar lenguaje imperativo con este modelo.
    \item La metodología orientada a objetos es una técnica o modelo, no depende del lenguaje.
    \item No es correcto suponer que este modelo sólo es útil con lenguajes orientados a objetos.
    \item Es un modelo versátil aplicable a pequeños y grandes proyectos independientemente del lenguaje.
    \item Según el documento (página 34): "Como mencionamos, es un modelo muy versátil, y por ser uno de los últimos en aparecer, aprendió mucho de los anteriores".
\end{enumerate}
\end{framed}

\subsection{Enumere el ciclo de vida y los pasos que seguiría, si debiese desarrollar una aplicación que monitoree el estado de las redes de una empresa.}

\begin{framed}
\noindent\textbf{Modificaciones sugeridas del documento:}
\begin{enumerate}
    \item Etapas iniciales: Expresión de necesidades, Especificaciones, Análisis y Diseño del sistema.
    \item Etapas finales: Implementación, Debugging, Validación y Evolución/Mantenimiento.
    \item Modelo recomendado según documento: Ciclo de Vida Orientado a Objetos.
    \item Justificación: Adecuado para programas de monitoreo de procesos y grandes sistemas de transacciones.
    \item Referencia documento (páginas 18-19 y 34): Detalla las 8 etapas generales y la aplicabilidad para monitoreo.
\end{enumerate}
\end{framed}

\subsection{Realice una lista de requerimientos hipotéticos para una aplicación que deba ejecutar archivos de música, pida la misma lista a un usuario no programador y compare las listas. ¿Qué enfoques encuentra en cada lista?}

\begin{framed}
\noindent\textbf{Modificaciones sugeridas del documento:}
\begin{enumerate}
    \item Objetivo de la etapa: Reflejar requerimientos y funcionalidades (qué, y no cómo, se va a implementar).
    \item Enfoque Programador (técnico): Formatos, Codecs, Gestión de memoria, API de audio y optimización.
    \item Enfoque Usuario no programador (funcional): Reproducir, Pausar, Volumen, Listas y facilidad de uso.
    \item Diferencia clave: El cliente o usuario sólo impartía especificaciones muy generales del producto final.
    \item Nota: Esta actividad es práctica y requiere trabajo externo con usuarios reales para comparar listas.
\end{enumerate}
\end{framed}

\subsection{A modo de encuesta, pregunte a sus colegas programadores, quién y ha utilizados un ciclo de vida. Indague sobre los resultados obtenidos.}

\begin{framed}
\noindent\textbf{Modificaciones sugeridas del documento:}
\begin{enumerate}
    \item Esta pregunta es una actividad de investigación externa con colegas programadores.
    \item Factores que influyen en la elección: Complejidad, Tiempo disponible y entregas parciales.
    \item Otros factores: Comunicación entre equipo y usuario, y certeza de los requerimientos.
    \item Resultados esperados: Diferentes programadores usan diferentes modelos según el tipo de proyecto.
    \item Según el documento (página 20): "Un consejo: la mejor metodología, es la que mejor conocemos".
\end{enumerate}
\end{framed}

\end{document}
